% ============================================================================
% Preámbulo: tipografía, color, cajas y estilos de diagramas
% ============================================================================
\usepackage[T1]{fontenc}
\usepackage[utf8]{inputenc}
\usepackage[spanish,es-noshorthands,es-tabla,es-nolists,es-nodecimaldot]{babel}
\usepackage[sc]{mathpazo}
\linespread{1.06}
\usepackage[scaled=0.90]{helvet}
\renewcommand\ttdefault{lmtt}
\usepackage{microtype}
\usepackage[letterpaper,textwidth=6.5in,textheight=8.9in,top=1.0in,headsep=0.28in,footskip=0.42in]{geometry}

\usepackage{amsmath,amssymb,mathtools,bm}
\usepackage{graphicx}
\usepackage[table]{xcolor}
\usepackage{tikz}
\usetikzlibrary{babel,positioning,arrows.meta,calc,fit,backgrounds,shapes.geometric,
  shapes.misc,decorations.pathreplacing,patterns,matrix,shadows}
\usepackage[most]{tcolorbox}
\usepackage{booktabs,tabularx,array,multirow,makecell}
\usepackage{enumitem}
\usepackage{caption}
\usepackage{subcaption}
\usepackage{float}
\usepackage{fancyhdr}
\usepackage{titlesec}
\usepackage{titletoc}
\usepackage{pifont}
\usepackage{listings}
\usepackage[round,sort]{natbib}
\usepackage{xurl}
\usepackage[hidelinks]{hyperref}
\usepackage{ragged2e}

% ---------------------------------------------------------------------------
% Paleta (misma que las figuras; orden categórico validado para daltonismo)
% ---------------------------------------------------------------------------
\definecolor{qazul}{HTML}{2A78D6}
\definecolor{qazuloscuro}{HTML}{104281}
\definecolor{qnaranja}{HTML}{EB6834}
\definecolor{qaqua}{HTML}{1BAF7A}
\definecolor{qamarillo}{HTML}{EDA100}
\definecolor{qvioleta}{HTML}{4A3AA7}
\definecolor{tinta}{HTML}{0B0B0B}
\definecolor{tinta2}{HTML}{52514E}
\definecolor{tenue}{HTML}{898781}
\definecolor{rejilla}{HTML}{E1E0D9}
\definecolor{eje}{HTML}{C3C2B7}
\definecolor{grisclaro}{HTML}{F0EFEC}
\definecolor{grisfondo}{HTML}{F7F7F5}
\definecolor{bueno}{HTML}{0CA30C}
\definecolor{alerta}{HTML}{FAB219}
\definecolor{serio}{HTML}{EC835A}
\definecolor{critico}{HTML}{D03B3B}
\definecolor{azulfondo}{HTML}{EEF4FC}
\definecolor{naranjafondo}{HTML}{FDF1EB}
\definecolor{aquafondo}{HTML}{E8F7F1}
\definecolor{alertafondo}{HTML}{FFF7E6}
\definecolor{buenofondo}{HTML}{EAF6EA}

\hypersetup{colorlinks=true,linkcolor=qazuloscuro,citecolor=qazuloscuro,urlcolor=qazuloscuro,
  pdftitle={QuantileFlow: seguimiento de distribuciones implícitas y gestión diaria de posiciones},
  pdfsubject={Propuesta técnica (versión de investigación)},pdfauthor={QuantileFlow},
  pdfkeywords={SSVI, Breeden-Litzenberger, Wasserstein, Kalman, BOCPD, CVaR, conformal}}

% ---------------------------------------------------------------------------
% Títulos, encabezados e índice
% ---------------------------------------------------------------------------
\setcounter{secnumdepth}{3}
\setcounter{tocdepth}{2}
\titleformat{\section}{\sffamily\LARGE\bfseries\color{tinta}}
  {\color{qazul}\thesection}{0.55em}{}[\vspace{3pt}{\color{eje}\titlerule[0.6pt]}]
\titlespacing*{\section}{0pt}{0pt}{12pt}
\titleformat{\subsection}{\sffamily\large\bfseries\color{tinta}}{\color{qazul}\thesubsection}{0.6em}{}
\titlespacing*{\subsection}{0pt}{16pt}{6pt}
\titleformat{\subsubsection}{\sffamily\normalsize\bfseries\color{tinta}}{\color{qazul}\thesubsubsection}{0.6em}{}
\titlespacing*{\subsubsection}{0pt}{11pt}{4pt}
\titleformat{\paragraph}[runin]{\sffamily\normalsize\bfseries\color{tinta}}{}{0pt}{}[.]

\titlecontents{section}[1.8em]{\addvspace{2.5pt}\sffamily\bfseries}
  {\contentslabel[\color{qazul}\thecontentslabel]{1.8em}}{}{\hfill\contentspage}
\titlecontents{subsection}[4.2em]{\small}
  {\contentslabel[\color{tinta2}\thecontentslabel]{2.4em}}{}{\titlerule*[0.6pc]{.}\contentspage}

\pagestyle{fancy}
\fancyhf{}
\renewcommand{\headrulewidth}{0.4pt}
\renewcommand{\headrule}{{\color{eje}\hrule width\headwidth height\headrulewidth}}
\fancyhead[L]{\sffamily\scriptsize\color{tinta2}QuantileFlow \textperiodcentered{} Propuesta técnica \textperiodcentered{} Versión de investigación}
\fancyhead[R]{\sffamily\scriptsize\color{tinta2}\nouppercase{\leftmark}}
\fancyfoot[C]{\sffamily\footnotesize\color{tinta2}\thepage}
\renewcommand{\sectionmark}[1]{\markboth{\thesection\quad #1}{}}
\fancypagestyle{plain}{\fancyhf{}\renewcommand{\headrulewidth}{0pt}\fancyfoot[C]{\sffamily\footnotesize\color{tinta2}\thepage}}

% ---------------------------------------------------------------------------
% Pies de figura, listas y tablas
% ---------------------------------------------------------------------------
\captionsetup{font={small,sf},labelfont={bf,color=qazuloscuro},labelsep=period,
  justification=justified,singlelinecheck=false,skip=6pt}
\captionsetup[table]{position=top}
\setlist{itemsep=2pt,topsep=4pt,parsep=0pt,leftmargin=1.4em}
\setlist[itemize,1]{label={\color{qazul}\small\textbullet}}
\setlist[itemize,2]{label={\color{tenue}\textendash}}
\setlist[enumerate,1]{label={\sffamily\bfseries\color{qazul}\arabic*.}}
\renewcommand{\arraystretch}{1.2}
\setlength{\aboverulesep}{0.3ex}
\setlength{\belowrulesep}{0.5ex}
\newcolumntype{Y}{>{\RaggedRight\arraybackslash}X}
\newcommand{\encab}[1]{\textsf{\bfseries\footnotesize #1}}

\setlength{\parindent}{0pt}
\setlength{\parskip}{5pt plus 1pt}
\setlength{\emergencystretch}{2em}
\raggedbottom
\renewcommand{\floatpagefraction}{0.75}
\renewcommand{\topfraction}{0.9}
\renewcommand{\textfraction}{0.08}

% ---------------------------------------------------------------------------
% Iconos
% ---------------------------------------------------------------------------
\newcommand{\iconoidea}{\tikz[baseline=-0.7ex]\node[circle,fill=qazul,text=white,
  font=\sffamily\bfseries\fontsize{6}{6}\selectfont,inner sep=0pt,minimum size=3mm]{i};}
\newcommand{\iconoalerta}{\tikz[baseline=-0.35ex]{\fill[alerta,rounded corners=0.3pt]
  (0,0)--(0.34,0)--(0.17,0.30)--cycle;\node[font=\sffamily\bfseries\fontsize{5.5}{5.5}\selectfont,
  text=tinta] at (0.17,0.105){!};}}
\newcommand{\iconocheck}{\tikz[baseline=-0.7ex]\node[circle,fill=bueno,text=white,
  font=\fontsize{6}{6}\selectfont,inner sep=0pt,minimum size=3mm]{\ding{51}};}
\newcommand{\iconopendiente}{\tikz[baseline=-0.7ex]\node[circle,draw=tenue,line width=0.6pt,
  text=tinta2,font=\sffamily\bfseries\fontsize{6}{6}\selectfont,inner sep=0pt,minimum size=3mm]{?};}
\newcommand{\icononota}{\tikz[baseline=-0.7ex]\node[rectangle,rounded corners=0.6pt,fill=qvioleta,
  text=white,font=\sffamily\bfseries\fontsize{6}{6}\selectfont,inner sep=0pt,minimum size=3mm]{N};}
\newcommand{\si}{\textcolor{bueno}{\ding{51}}}
\newcommand{\no}{\textcolor{critico}{\ding{55}}}

% ---------------------------------------------------------------------------
% Cajas temáticas
% ---------------------------------------------------------------------------
\tcbset{
  qfbase/.style={enhanced,breakable,boxrule=0pt,arc=0pt,outer arc=0pt,left=9pt,right=9pt,
    top=6pt,bottom=6pt,fontupper=\small,before skip=9pt,after skip=9pt,
    fonttitle=\sffamily\bfseries\small,coltitle=tinta,attach title to upper={\par\smallskip}},
}
\newtcolorbox{ideaclave}[1][Idea clave]{qfbase,colback=azulfondo,
  borderline west={2.2pt}{0pt}{qazul},title={\iconoidea\hspace{0.45em}#1}}
\newtcolorbox{ideaclavefinal}[1][Idea clave]{qfbase,unbreakable,colback=azulfondo,
  borderline west={2.2pt}{0pt}{qazul},title={\iconoidea\hspace{0.45em}#1}}
\newtcolorbox{noconcluir}[1][Qué no concluir]{qfbase,colback=alertafondo,
  borderline west={2.2pt}{0pt}{alerta},title={\iconoalerta\hspace{0.45em}#1}}
\newtcolorbox{criterio}[1][Criterio para avanzar]{qfbase,colback=buenofondo,
  borderline west={2.2pt}{0pt}{bueno},title={\iconocheck\hspace{0.45em}#1}}
\newtcolorbox{pendiente}[1][Pendiente de definir]{qfbase,colback=grisfondo,
  borderline west={2.2pt}{0pt}{tenue},title={\iconopendiente\hspace{0.45em}#1}}
\newtcolorbox{notatecnica}[1][Nota técnica]{qfbase,colback=grisfondo,
  borderline west={2.2pt}{0pt}{qvioleta},title={\icononota\hspace{0.45em}#1}}
\newtcolorbox{algoritmo}[1]{qfbase,colback=white,colframe=eje,boxrule=0.5pt,
  borderline west={0pt}{0pt}{white},title={#1},fontupper=\small}

% Ficha de etapa: entradas, proceso, salidas y controles
\newcommand{\ficha}[4]{%
  \begin{tcolorbox}[enhanced,boxrule=0pt,arc=1.5pt,colback=grisfondo,left=7pt,right=7pt,top=6pt,
    bottom=5pt,before skip=4pt,after skip=12pt]
  \small
  \begin{tabularx}{\linewidth}{@{}>{\RaggedRight}X@{\hspace{10pt}}>{\RaggedRight}X@{\hspace{10pt}}>{\RaggedRight}X@{\hspace{10pt}}>{\RaggedRight\arraybackslash}X@{}}
  {\sffamily\bfseries\footnotesize\color{qazul}ENTRADAS} &
  {\sffamily\bfseries\footnotesize\color{qazul}PROCESO} &
  {\sffamily\bfseries\footnotesize\color{qazul}SALIDAS} &
  {\sffamily\bfseries\footnotesize\color{qazul}CONTROLES}\\[2pt]
  \footnotesize #1 & \footnotesize #2 & \footnotesize #3 & \footnotesize #4
  \end{tabularx}
  \end{tcolorbox}}

% ---------------------------------------------------------------------------
% Estilos TikZ
% ---------------------------------------------------------------------------
\tikzset{
  >={Stealth[length=2.1mm,width=1.5mm]},
  base/.style={font=\sffamily\scriptsize,text=tinta,align=center},
  caja/.style={base,draw=#1,fill=#1!7,rounded corners=2pt,line width=0.6pt,inner sep=4pt,
    minimum height=7mm},
  caja/.default=tenue,
  cajallena/.style={base,draw=none,fill=#1,text=white,rounded corners=2pt,inner sep=4pt,
    minimum height=7mm,font=\sffamily\scriptsize\bfseries},
  flecha/.style={->,line width=0.6pt,draw=tinta2},
  flechagruesa/.style={->,line width=1.1pt,draw=tinta2},
  guia/.style={line width=0.5pt,draw=eje},
  nota/.style={font=\sffamily\fontsize{6.4}{7.6}\selectfont,text=tinta2,align=center},
  titulo/.style={font=\sffamily\scriptsize\bfseries,text=tinta},
  capa/.style={rounded corners=3pt,fill=#1!5,draw=#1!45,line width=0.5pt},
}

% Listados (registro de corridas)
\lstdefinestyle{registro}{basicstyle=\ttfamily\scriptsize,columns=fullflexible,keepspaces=true,
  frame=none,backgroundcolor=\color{grisfondo},xleftmargin=6pt,xrightmargin=6pt,aboveskip=6pt,
  belowskip=6pt,showstringspaces=false,stringstyle=\color{qazuloscuro},
  literate={á}{{\'a}}1 {é}{{\'e}}1 {í}{{\'i}}1 {ó}{{\'o}}1 {ú}{{\'u}}1 {ñ}{{\~n}}1}

% ---------------------------------------------------------------------------
% Notación
% ---------------------------------------------------------------------------
\newcommand{\Q}{\mathbb{Q}}
\renewcommand{\P}{\mathbb{P}}
\newcommand{\E}{\mathbb{E}}
\newcommand{\dd}{\,\mathrm{d}}
\newcommand{\CVaR}{\operatorname{CVaR}}
\newcommand{\VaR}{\operatorname{VaR}}
\newcommand{\sintetico}{\textsf{\scriptsize\color{tenue}[datos sintéticos]}}
\newcommand{\grafica}[4][0.66]{%
  \begin{figure}[H]\centering
  \includegraphics[width=#1\linewidth]{figuras/#2}
  \caption{#3}\label{#4}
  \end{figure}}
\newcommand{\figref}[1]{figura~\ref{#1}}
\newcommand{\Figref}[1]{Figura~\ref{#1}}
\newcommand{\secref}[1]{sección~\ref{#1}}
\newcommand{\tabref}[1]{tabla~\ref{#1}}
